%-*-latex-*-
%-*-latex-*-
\newcommand\COURSE{ciss350}
\newcommand\ASSESSMENT{q10604}
\newcommand\ASSESSMENTTYPE{Quiz}
\newcommand\POINTS{\textwhite{xxx/xxx}}

\input{myquizpreamble}
\input{yliow}
\input{\COURSE}
%-*-latex-*-

%-*-latex-*-

%-*-latex-*-
\renewcommand\TITLE{\ASSESSMENTTYPE \ \ASSESSMENT}

\renewcommand\EMAIL{}
\renewcommand\AUTHOR{eagraham1@cougars.ccis.edu}

\textwidth=6in
\begin{document}
\topmatterthree

%------------------------------------------------------------------------------
\nextq
Infix to Polish notation:
\\
Convert the following in infix notation \verb!1 + (2 + 3) * 4 + 5! to
Polish notation.
If what is given has an error, write ERROR.
Note that you are converting the expression, you are not
evaluating the expression.
\\
\ANSWER
\begin{answercode}
+ + 1 * + 2 3 4 5
\end{answercode}

%------------------------------------------------------------------------------
\nextq
Polish to infix notation:
\\
Convert the following in Polish notation \verb!+ / 8 2 - 4 * 7 2!
to infix notation.
If what is given has an error, write ERROR.
Note that you are converting the expression, you are not
evaluating the expression.
\\
\ANSWER
\begin{answercode}
8 / 2 + (4 - 7 * 2)
\end{answercode}

%------------------------------------------------------------------------------
\nextq
Infix to Reverse Polish notation:
\\
Convert the following in infix notation \verb!1 / (2 + 3 * 4) * 5! to
Reverse Polish notation.
If what is given has an error, write ERROR.
Note that you are converting the expression, you are not
evaluating the expression.
\\
\ANSWER
\begin{answercode}
1 2 3 4 * + / 5 *
\end{answercode}

%------------------------------------------------------------------------------
\nextq
Reverse Polish to infix notation:
\\
Convert the following in Reverse Polish notation \verb!2 3 4 * 5 * +!
to infix notation.
If what is given has an error, write ERROR.
Note that you are converting the expression, you are not
evaluating the expression.
\\
\ANSWER
\begin{answercode}
2 + (3 * 4 * 5)
\end{answercode}

%------------------------------------------------------------------------------
\nextq
Using five \verb!1!s and as many \verb!+!s as you like,
write a mathematical expression in Reverse Polish notation
that will have the largest possible
stack during the evaluation of the expression.
\\
\ANSWER
\begin{answercode}
1 1 1 1 1 + + + +
\end{answercode}

%------------------------------------------------------------------------------
\nextq
Using five \verb!1!s and as many \verb!+! as you like,
write a mathematical expression in Reverse Polish notation
that will have the smallest possible
stack during the evaluation of the expression.
\\
\ANSWER
\begin{answercode}
1 1 + 1 + 1 + 1 +
\end{answercode}

%------------------------------------------------------------------------------
\nextq
Complete the following function that returns \verb!true! if the vector
of numeric and operator strings is a valid RPN expression.
For instance \verb!{"1", "2", "+"}! is valid and
\verb!{"1", "+", "2"}! is not valid.
Note that you are not evaluating the expression.
The following functions are provided:
\begin{Verbatim}[fontsize=\footnotesize,frame=single]
bool is_numeric(const std::string & s);
bool is_operator(const std::string & s);
\end{Verbatim}
\ANSWER
\begin{answercode}
#include <iostream>
#include <string>
#include <vector>
#include <stack>

// Returns true if string is numeric.
// Assume s.size() > 0 and no sign ('+' or '-')
bool is_numeric(const std::string & s)
{
    if (s.size() == 0) return false;
    for (unsigned int i = 0; i < s.size(); ++i)
    {
        if (s[i] < '0' || '9' < s[i])
        {
            return false;
        }
    }
    return true;
}

// Returns true if string is an operator.
// Assume operator is +, -, *, /, %
bool is_operator(const std::string & s)
{
    if (s.size() != 1) return false;
    return (s[0] == '+' || s[0] == '-' || s[0] == '*' || s[0] == '/' ||
            s[0] == '%');
}

bool is_rpn(const std::vector< std::string > & v)
{
  int count = 0;
  
  for (unsigned int i = 0; i < v.size(); ++i)
  {
    std::string t = v[i];

    if(is_numeric(t))
    {
      count++;
    }
    else if (is_operator(t))
    {
      if (count < 2)
          return false;
      count--;
    }
    else
    {
      return false;
    }
  }
  return count == 1;
}

int main()
{
    std::vector< std::string > v {"1", "2", "+"};
    std::cout << is_rpn(v) << '\n';
    return 0;
}
\end{answercode}
(Hint: Do you really need a stack?)

%------------------------------------------------------------------------------
%-*-latex-*-
\newpage
\input{instructions}
\end{document}

